% Pursuing Stacks (cotes 134-2 to 134-8): archive cote 134-4, site page 77 (PDF page 78), typescript page 165.
% typescript. Draft compared: ps3.tex:3931-4002.
% AI-transcribed by Claude Opus 5.5 (claude-opus-5-5, adaptive thinking, effort high), Message Batches API, batch ps2 (msgbatch_01NVwr1Yuo8qJf5rykeGBXPK), 2026-11-12; not yet checked by eye.
\scanpage{134-4/77}
\mspage{165}

\uline{Definition}  A subset \underline{W} of F\unsure{l}((Cat)) is called a \uline{basic} \uline{localizer} ,  if it
satisfies the conditions L 1),L 3), L 4) above (saturation, final object
and localization axioms),  and hence also the homotopy axiom L 2).

These conditions are enough,  I quickly checked this night,  in or-%
der to validify all results developped so far on test categories, weak
test categories,  strict test categories,  weak test functors and test
functors \inserted{with values in (Cat)} (cf notably the review in par. 44, page 79 - 88),   provided
in the case of test functors we restrict to the case of loc. cit. when
each of the categories i(a) has a final object. \inserted{All} \over{T}{t}his I believe is justi-%
fication enough for the definition above.

As for the fibration axiom\xout{\unsure{(}} L 5),  this we have seen to be needed
(at least in the approach we got so far) for handling test functors
i: $A \longrightarrow$ (Cat),  while no longer assuming the categories i(a) to have final
objects\hand{,} \hand{(}which was felt to be a significant generalization to carry
through,  in view of being able subsequently to replace (Cat) by more
general modelizers\hand{)}. While still in the nature of a stability requirement\hand{,}
this fibration axiom looks to me a great deal stronger than the other
axioms.  Clearly L 5),  together with the very weak \inserted{``interval axiom''} L 3' ($\hand{\Delta}_1 \hand{\to}$ e is in
\hand{\unsure{$\vee$}}\underline{W} )  implies the homotopy axiom.  It can be seen too that when joined
\margin{\unsure{this plus the fact that L5+L3'} +L1 \unsure{implies} \unsure{W=W$^0$} ! \unsure{If this not known, we need L5* + this condition} \illegible{2} \unsure{cf p. 30} \illegible{1} \unsure{or}}
with L 1),  it implies the localization axiom L 4 (using the standart
device of a mapping-cone for a functor ...)\hand{*}.  Thus a basic \over{\unsure{f}}{l}ocalizer
satisfying the fibration axiom L 5)  can be viewed also as a \underline{W}  satisfy-%
ing the conditions

L 1 (saturation),  L 3') (interval axiom) , L 5) (fibration axiom).

In the next section,  after we will have stated the two key-\deleted{statements}\inserted{facts}
about weak test functors and test functors,  which both make use of L 5),
we'll presumably\hand{,} for the rest of the work ahead towards canonical modeli-%
zers,  assume \deleted{that} \inserted{the fibration} condition on the basic \over{\unsure{f}}{l}ocalizer \underline{W} .

There  are some other properties of weak equivalence $W_{(Cat)}$ and it's
manifold variants in terms of cohomology,  which have not been listed
yet,  and which surely will turn up still sooner or later.  Maybe it's
too soon to line them up in a definite order,  as their \xout{exact} significan-%
ce is still somewhat vague  and needs closer scrutiny.  I'll just list
those which come to my mind,  in a provisional order. \deleted{\illegible{2}}

\hand{L} a)  O-admissibility of \underline{W} , namely $f \hand{\in} \underline{W}$  implies $\hand{\pi}_0$(f) bijective.

This condition,  together with the homotopy axiom L 2),  will imply

\margin{77}
